\documentclass[11pt]{article}

\usepackage[margin=25mm]{geometry}
\usepackage[T1]{fontenc}
\usepackage[utf8]{inputenc}
\usepackage{amsmath,amssymb}
\usepackage{graphicx}
\usepackage{booktabs,tabularx,array}
\usepackage{lineno}
\usepackage[numbers,sort&compress]{natbib}
\usepackage[hidelinks]{hyperref}
\usepackage{xcolor}
\setlength{\emergencystretch}{3em}

\graphicspath{{../outputs/final/publication/figures/}}
\newcommand{\Terafab}{Terafab}
\newcommand{\wrated}{\mathrm{W_{rated}}}
\newcommand{\code}[1]{\texttt{\detokenize{#1}}}
\newcommand{\limitation}[1]{\textit{Boundary: #1}}

\title{From terawatt-scale semiconductor output to grid burden:\\
A prospective constraint assessment of the announced Terafab project}
\author{Khalid M. Saqr\\
\small KNOWDYN Ltd, United Kingdom}
\date{}

\begin{document}
\linenumbers
\maketitle

\begin{abstract}
Terafab has publicly described semiconductor output of 1 TW/year, but the
dimension of that target and the associated facility requirements are not
defined. This study asks whether currently available public evidence
demonstrates a realistic project and what the admitted interpretation would
imply for United States electricity security through 2050. A prospective,
evidence-gated model translates one terawatt of aggregate rated power embodied
in good compute devices manufactured per year into wafer throughput, facility
electricity, coincident load, heat rejection and water withdrawal. It then
tests 337 predeclared deployment--supply--stress trajectories against tri-state
manufacturing, infrastructure and ERCOT adequacy gates, with correlated
uncertainty propagated by scrambled Sobol sampling. Under the central
full-target interpretation, the model requires 68.29 million 300-mm wafer
starts/year, 135.17 TWh/year, 16.58 GW coincident load and heat rejection, and
568.72 million m$^3$/year of external water withdrawal. The modeled electricity
is 2.64\% of the EIA baseline US total in 2037 and 2.26\% in 2050. In the locked
correlated uncertainty case, median electricity is 125.11 TWh/year, whereas
marginal p95 peak and water requirements reach 91.54 GW and 1.571 billion
m$^3$/year. All 336 nonzero trajectories remain indeterminate because public
evidence does not establish target semantics, phase manufacturing capacity,
utility capacity, executable interconnection or probabilistic ERCOT adequacy.
Thus public evidence does not presently demonstrate project realism; this is
not proof of infeasibility. The contribution is a reproducible method for
separating resource implications from verified project facts.
\end{abstract}

\noindent\textbf{Keywords:} semiconductor manufacturing; electricity security;
large loads; ERCOT; uncertainty; energy--water nexus; prospective assessment

\section{Introduction}
\label{sec:introduction}

The announced \Terafab{} concept joins two energy-system questions that are
usually treated separately: how a claimed semiconductor production scale maps
to physical manufacturing throughput, and how the resulting industrial load
interacts with an electricity system already absorbing rapidly growing large
loads. The public website states ``Terafab output 1 TW/Year''
\cite{Terafab2026}; Texas incentive applications identify four prospective
phases and schedule information \cite{JETI2026}. Neither source supplies a
dimensional definition of the target, phase wafer capacity, facility demand,
water entitlement, or an executable grid interconnection. Equating the public
target directly with factory electrical load would therefore be dimensionally
and evidentially invalid.

Semiconductor production is a coupled electricity--thermal--water process.
Cleanroom air handling, process cooling and ultrapure-water systems can be
material components of facility demand, while their intensities vary greatly
across sites and process generations \cite{Hu2019,Mathew2010,Tschudi2003}.
Manufacturer reports provide useful scale comparators, but company aggregates
span multiple sites, nodes and products \cite{TSMCAnnual2024,TSMCSustainability2024}.
Water-demand estimates are similarly heterogeneous and cannot establish a
site-specific entitlement \cite{Sandhu2025}. A study of an announced facility
must therefore distinguish public project evidence, industry comparators,
scenario assumptions and derived results.

The regional context adds a second identification problem. ERCOT planning
reports and effective-load-carrying-capability studies support accredited-
capacity screening through their stated horizons \cite{ERCOTCDR2025,ERCOTELCC2024},
but reserve margin alone is not the Texas probabilistic reliability standard
\cite{PUCT25508}. National EIA projections contextualize annual energy scale
but do not determine ERCOT adequacy \cite{EIAAEO2026}. NERC has separately
identified the planning relevance and uncertainty of emerging large loads
\cite{NERCLTRA2025,NERCLargeLoads2025}.

This paper asks: (i) what manufacturing and facility-resource requirements
follow from a transparent interpretation of the target; (ii) when those
requirements arise under predeclared realization pathways; (iii) what they
imply for ERCOT capacity screening and US annual electricity; and (iv) whether
the admitted public evidence is sufficient for a realism classification. Its
novelty is not a prediction of the project. It is a dimensionally consistent,
evidence-gated production-to-energy framework that refuses to convert unknown
capacity or reliability inputs into favorable or unfavorable findings.

\section{Methods}
\label{sec:methods}

\subsection{Evidence contract, boundary and target semantics}

All final calculations use frozen, hashed extracts dated 1 August 2026; the
reproduction path performs no live web retrieval. Sources were classified as
official project publications, public filings, manufacturer disclosures,
peer-reviewed or government-laboratory benchmarks, official system
projections, regulations and technical studies. Reported and filed claims were
admitted only for their stated roles. Restricted, private, proprietary and
unpublished project material was excluded.

The primary conditional interpretation is

\begin{quote}
one terawatt of aggregate rated electrical power embodied in good compute
devices manufactured per calendar year.
\end{quote}

The annual target has unit $\wrated$/year. It is product output, not site
electrical power. An alternative public semantic branch remains explicit but
non-numeric because the public wording does not identify another calculable
quantity. Figure~\ref{fig:architecture} shows the claim-to-consequence chain,
and Table~\ref{tab:claims} records the excluded interpretations.

\begin{figure*}[htbp]
\centering
\includegraphics[width=0.94\textwidth]{figure_01.pdf}
\caption{Claim-to-consequence architecture. Product rated output is translated
through device and wafer throughput before facility electricity, heat and water
are calculated; the public output wording is never assigned as factory load.}
\label{fig:architecture}
\end{figure*}

% Compact representations of the six prespecified tables. The complete,
% machine-readable versions are generated under outputs/final/publication/tables.

\begin{table*}[htbp]
\centering
\caption{Public claims and admitted interpretations. “Reported” and “filed”
describe evidence provenance, not verification of project performance.}
\label{tab:claims}
\small
\begin{tabularx}{\textwidth}{p{0.14\textwidth}p{0.13\textwidth}p{0.13\textwidth}X X}
\toprule
Public wording & Source/status & Admitted interpretation & Excluded interpretation & Consequence \\
\midrule
``Terafab output 1 TW/Year'' & Public website; reported & Aggregate rated electrical power embodied in good compute devices manufactured per calendar year & Factory electrical demand; deployed-device consumption; computing performance & Numerical conditional target; semantic gate remains indeterminate \\
Four Texas phases with first-operation or incentive years 2029, 2032, 2035 and 2037 & JETI applications; filed claims & Scenario schedule only & Completed capacity, utility entitlement or executable interconnection & Schedule branches; phase output scale remains unknown \\
\bottomrule
\end{tabularx}
\end{table*}

\begin{table*}[htbp]
\centering
\caption{Principal variables and governing equations. Complete definitions,
units, temporal bases, evidence classes and module ownership are in the generated
Table 2.}
\label{tab:equations}
\small
\begin{tabularx}{\textwidth}{p{0.11\textwidth}p{0.23\textwidth}p{0.18\textwidth}X}
\toprule
Symbol & Definition & Unit & Governing relation \\
\midrule
$D$ & Gross dies per 300-mm wafer & die wafer$^{-1}$ & $D=\pi d^2/(4A)-\pi d/\sqrt{2A}$ \\
$W_y$ & Annual wafer starts & wafer yr$^{-1}$ & $W_y=C_{rated}N_d/(DYP_m)$ \\
$E_f$ & Total facility electricity & MWh yr$^{-1}$ & Throughput times admitted total-facility intensity \\
$P_f$ & Coincident facility peak & MW & $P_f=E_f/(8760\lambda)$ \\
$Q_{rej}$ & Peak heat rejection & MW$_{th}$ & $Q_{rej}=P_f$ under the stated steady balance \\
$V_w$ & External water withdrawal & m$^3$ yr$^{-1}$ & Gross demand less declared reuse; no double credit \\
$P_g$ & Grid coincident peak & MW & $P_g=\max(0,P_f-C_{onsite,accr})$ \\
$RM$ & Reserve-margin screen & fraction & $(C_{accr}-P_{peak})/P_{peak}$ \\
$s^E$ & US annual-electricity share & fraction & $E_f/E^{US,0}$ \\
\bottomrule
\end{tabularx}
\end{table*}

\begin{table*}[htbp]
\centering
\caption{Evidence and validation dataset. These comparators define ranges or
diagnostics; they do not constitute Terafab measurements.}
\label{tab:evidence}
\small
\begin{tabularx}{\textwidth}{p{0.18\textwidth}p{0.16\textwidth}p{0.18\textwidth}X p{0.17\textwidth}}
\toprule
Source & Context & Metric role & Use & Limitation \\
\midrule
TSMC annual and sustainability reports & Mixed corporate fabs & Wafer, electricity and water scale & Comparator/range & Company aggregate, not a project analogue \\
Hu et al. \cite{Hu2019} & Taiwanese high-tech fab & Facility/component energy & Calibration component & Single fab and non-project process \\
Mathew et al. \cite{Mathew2010} & Fabs21 cleanrooms & Energy heterogeneity & Range diagnostic & Historical mixed cleanrooms \\
Tschudi and Xu \cite{Tschudi2003} & Fourteen cleanrooms & HVAC energy & Cooling reference & Historical technology and climates \\
Sandhu et al. \cite{Sandhu2025} & 300-mm semiconductor review & Water intensity & Water reference & Review estimate, not a site measurement \\
\bottomrule
\end{tabularx}
\end{table*}

\begin{table*}[htbp]
\centering
\caption{Predeclared scenario matrix. The Cartesian experiment contains one
no-build counterfactual and 336 nonzero trajectories.}
\label{tab:matrix}
\small
\begin{tabularx}{\textwidth}{p{0.17\textwidth}p{0.16\textwidth}p{0.17\textwidth}X}
\toprule
Axis & Levels & Numerical status & Interpretation boundary \\
\midrule
Target scale & no build; research fab; initial large scale; full announced target & 0; unknown; unknown; conditional 1 TW$_{rated}$/yr & Labels do not create capacity evidence \\
Realization & accelerated; reference; delayed; not realized by 2050 & Annual fractions, 2026--2050 & No pathway is official or predictive \\
Supply & grid dominant; firm onsite; renewable--storage--grid; firm low-carbon hybrid & Conditional capacities and accreditation & No portfolio is claimed as the project design \\
Stress & normal plus six schedule, weather, yield, water/cooling and outage stresses & One declared stress per branch & Stress failures are not ordinary-case project facts \\
Counterfactual & ERCOT protocol and SB6 through 2030; EIA baseline and high demand through 2050 & Official horizon only & No ERCOT extrapolation after 2030 \\
\bottomrule
\end{tabularx}
\end{table*}

\begin{table*}[htbp]
\centering
\caption{Principal conditional results. Facility quantities at full target are
identical across realization pathways; portfolios allocate the modeled peak
between accredited onsite supply and the grid. Every row remains indeterminate.}
\label{tab:principal}
\small
\resizebox{\textwidth}{!}{%
\begin{tabular}{llrrrrrrl}
\toprule
Pathway/portfolio & Year & Target & Wafers & Electricity & Peak/heat & Withdrawal & Net grid peak & Class \\
 & & fraction & million yr$^{-1}$ & TWh yr$^{-1}$ & GW & million m$^3$ yr$^{-1}$ & GW & \\
\midrule
Accelerated/grid dominant & 2029 & 0.25 & 17.07 & 33.79 & 4.15 & 142.18 & 4.15 & indeterminate \\
Accelerated/grid dominant & 2037 & 1.00 & 68.29 & 135.17 & 16.58 & 568.72 & 16.58 & indeterminate \\
Reference/grid dominant & 2050 & 1.00 & 68.29 & 135.17 & 16.58 & 568.72 & 16.58 & indeterminate \\
Delayed/grid dominant & 2050 & 1.00 & 68.29 & 135.17 & 16.58 & 568.72 & 16.58 & indeterminate \\
Not realized/grid dominant & 2050 & 0.15 & 10.24 & 20.28 & 2.49 & 85.31 & 2.49 & indeterminate \\
Accelerated/firm onsite & 2050 & 1.00 & 68.29 & 135.17 & 16.58 & 568.72 & 3.32 & indeterminate \\
Accelerated/renewable--storage & 2050 & 1.00 & 68.29 & 135.17 & 16.58 & 568.72 & 5.47 & indeterminate \\
\bottomrule
\end{tabular}}
\end{table*}

\begin{table*}[htbp]
\centering
\caption{Verification, validation and reproduction status.}
\label{tab:verification}
\small
\begin{tabularx}{\textwidth}{p{0.21\textwidth}p{0.19\textwidth}p{0.13\textwidth}p{0.12\textwidth}X}
\toprule
Check & Metric & Result & Criterion & Interpretation \\
\midrule
Rated-output identity & Relative residual & 0 & $\leq10^{-8}$ & Passed on synthetic complete case \\
Facility energy balance & Relative residual & 0 & $\leq10^{-8}$ & Passed algebraic allocation identity \\
Water balance & Relative residual & 0 & $\leq10^{-8}$ & Passed when the admitted fraction is defined \\
Fixed-seed reproduction & Scenario hash & bitwise repeat & exact & Passed automated API test \\
Uncertainty convergence & Worst relative change & 0.002900 & $<0.01$ & Passed across independence and dependence stresses \\
Independent validation & Numeric holdouts & 0 & $\geq5$ & Insufficient; no empirical accuracy claim \\
CLI/notebook parity & Shared API & passed & $10^{-10}$ & Governing equations are not duplicated \\
Claim level & Model description & prospective constraint analysis & not validated twin & Passed \\
\bottomrule
\end{tabularx}
\end{table*}


\subsection{Production translation}

For a 300-mm wafer of diameter $d$ and die area $A$, the continuous edge-loss
approximation is

\begin{equation}
D=\frac{\pi d^2}{4A}-\frac{\pi d}{\sqrt{2A}},
\label{eq:dies}
\end{equation}

where $D$ is gross dies per wafer. Scribe-lane, defect and packaging losses are
not hidden in Eq.~\eqref{eq:dies}; they are represented through die area and
effective yield. If a compute module has rated power $P_m$, contains $N_d$ logic
dies and has effective yield $Y$, annual wafer starts are

\begin{equation}
W_y=\frac{C_{rated}N_d}{DYP_m}.
\label{eq:target}
\end{equation}

Reconstructing $C_{rated}$ from the calculated throughput is required to close
within a relative residual of $10^{-8}$. Figure~\ref{fig:surface} maps the full
target requirement across the admitted yield and module-power ranges. Because
no public project wafer capacity is available, it is a requirement surface,
not a pass/fail feasibility frontier.

\begin{figure}[htbp]
\centering
\includegraphics[width=\linewidth]{figure_02.pdf}
\caption{Full-target wafer-start requirement across effective yield and rated
module power. Admitted ranges and extrapolated regions are distinguished; no
project capacity boundary is inferred.}
\label{fig:surface}
\end{figure}

\subsection{Facility electricity, thermal and water balances}

Total annual facility electricity $E_f$ is throughput multiplied by an admitted
total-facility electricity intensity. The source already includes cooling, so
cooling cannot be added a second time. For non-cooling electricity $B$ and
coefficient of performance $COP$,

\begin{equation}
E_f=B+\frac{B}{COP},\qquad B=E_f\frac{COP}{COP+1}.
\label{eq:cooling}
\end{equation}

Coincident peak is derived from annual electricity and load factor $\lambda$,
$P_f=E_f/(8760\lambda)$. Under the declared annual steady balance, negligible
stored energy and no evidenced useful heat sink, the complete coincident
electrical load is reported as peak heat rejection. Recoverable heat is not
quantified without temperature-grade and sink evidence.

Water intensities are explicitly typed as either gross process demand before
reuse or external withdrawal after reuse. Reuse is subtracted only from the
former, preventing double credit. When consumption is defined, the steady
balance requires

\begin{equation}
V_{withdraw}=V_{consumed}+V_{discharged}
\label{eq:water}
\end{equation}

to a relative residual of $10^{-8}$. Table~\ref{tab:evidence} separates the
calibration comparators from unavailable independent project validation.
Figure~\ref{fig:benchmark} is consequently a benchmark diagnostic, not a
measured-versus-modeled validation result.

\begin{figure}[htbp]
\centering
\includegraphics[width=\linewidth]{figure_03.pdf}
\caption{Benchmark diagnostic and validation boundary. The same public
comparators that inform admitted ranges cannot serve as independent Terafab
holdouts; no bias, NRMSE, MAPE or coverage statistic is manufactured.}
\label{fig:benchmark}
\end{figure}

\subsection{Deployment, supply and electricity-security screening}

The predeclared matrix comprises one no-build counterfactual and 336 nonzero
trajectories: three nonzero target labels, four realization pathways, four
supply portfolios and seven stress conditions (Table~\ref{tab:matrix}). The
research-fab and initial-large-scale labels remain non-numeric because public
phase capacities are unknown. Numerical annual calculations cover 2026--2050
only for the full-target interpretation.

For hour $h$, grid load is conceptually $L^*_{t,h}=L^0_{t,h}+P_{grid,t,h}$.
In the annual screening model, accredited onsite capacity is subtracted only
after commissioning and availability rules are applied. Reserve margin is

\begin{equation}
RM_t=\frac{C_{accredited,t}-P_{peak,t}}{P_{peak,t}}.
\label{eq:rm}
\end{equation}

The regional adequacy gate can pass only when the applicable loss-of-load
expectation, maximum expected event duration and expected highest hourly
average load shed are all available and satisfy the Texas standard. Missing
probabilistic values produce \emph{indeterminate}; reserve margin cannot be
promoted to a reliability verdict. Official ERCOT counterfactuals are used only
through 2030, with no extrapolation. National scale is
$s_t^E=E_{f,t}/E_{t}^{US,0}$ against the EIA baseline and high-demand cases.

\subsection{Decision rules and uncertainty}

Ten gates cover target semantics, manufacturing throughput, facility energy,
thermal capacity, water capacity, schedule precedence, supply deliverability,
regional adequacy, national context and evidence sufficiency. Every gate is
\emph{pass}, \emph{fail} or \emph{indeterminate}. A pathway is conditionally
feasible only if all required gates pass. Explicit failures are retained, but
an indispensable unknown prevents an overall infeasibility finding. No weighted
energy-security index is constructed.

Ten manufacturing and facility parameters are propagated with a scrambled
Sobol design \cite{Owen1998}, fixed seed 20260801, 16,384 preliminary samples
and 32,768 confirmation samples. A Gaussian copula implements the predeclared
dependence structure and $\pm0.20$ correlation stresses. Convergence requires
strictly less than 1\% change in estimable medians and p95 quantities when the
sample count is doubled. Independent-input first- and total-order Sobol indices
\cite{Saltelli2010} are supplemented by a dependence-aware Shapley allocation
\cite{OwenPrieur2017}; 1,000 bootstrap replicates quantify ranking uncertainty.
Sensitivity is diagnostic for conditional facility peak and cannot substitute
for the locked but unidentifiable joint-feasibility outcome.

\subsection{Software implementation and reproducibility}

The source-available Python overlay calls a single typed kernel from both the
command line and notebook. It exports full scenario, annual and gate tables;
one CSV per figure; six paper tables; uncertainty and hypothesis JSON; and a
SHA-256 manifest. A synthetic complete case verifies algebraic identities but
is never admitted as a Terafab scenario. The complete command, environment and
artifact map are provided in the supplementary material.

OpenAI Codex assisted implementation of source code, automated checks,
deterministic Matplotlib visualization code and manuscript organization. It did
not retrieve live data in the final reproduction path, fabricate observations,
or alter generated numerical results. The author remains responsible for
source admission, scientific judgement and final verification; the formal
declaration appears before the references.

\section{Results}

\subsection{Central conditional requirement}
\label{sec:results-central}

In the accelerated pathway, modeled production reaches 25\% of the full target
in 2029, 50\% in 2032, 75\% in 2035 and 100\% in 2037
(Figure~\ref{fig:pathways}). Full target requires 68.2907 million wafer
starts/year (5.6909 million/month), 135.1736 TWh/year of facility electricity,
16.5809 GW coincident load, 16.5809 GW peak heat rejection and 568.7183 million
m$^3$/year external water withdrawal. These values are identical after full
target across the accelerated, reference and delayed demand pathways because
the same production interpretation and central facility parameters apply; the
pathways change timing rather than target magnitude. Table~\ref{tab:principal}
reports selected portfolio allocations.

\begin{figure*}[htbp]
\centering
\includegraphics[width=0.94\textwidth]{figure_04.pdf}
\caption{Conditional full-target realization pathways from 2026 to 2050.
Research-fab and initial-large-scale branches are not plotted numerically because
their public production capacities are unknown.}
\label{fig:pathways}
\end{figure*}

Figure~\ref{fig:decomposition} separates electricity, cooling and water
quantities without double-counting cooling electricity or water reuse. The heat
result is a first-law rejection requirement, not an identified useful-heat
resource. Water entitlement, consumption fraction and discharge infrastructure
remain project-specific unknowns outside the ordinary central case.

\begin{figure*}[htbp]
\centering
\includegraphics[width=0.94\textwidth]{figure_05.pdf}
\caption{Facility thermodynamic and water decomposition at principal
milestones. Recoverable heat is omitted because temperature-grade and sink
evidence are unavailable.}
\label{fig:decomposition}
\end{figure*}

\subsection{ERCOT screening and supply portfolios}
\label{sec:results-ercot}

At 25\% target in 2029--2030, the central grid-dominant requirement is 4.145 GW
coincident load and 33.793 TWh/year. Figure~\ref{fig:ercot} reports the two
frozen ERCOT counterfactuals only through 2030. Under the protocol-prescribed
case, the scenario's assumed accredited grid addition produces a positive
reserve-margin change; under the SB6 counterfactual, the sign can reverse. This
counterfactual dependence prevents interpretation of a positive screening
number as a project reliability benefit.

\begin{figure}[htbp]
\centering
\includegraphics[width=\linewidth]{figure_06.pdf}
\caption{ERCOT incremental burden through the last frozen official
counterfactual year. Reserve-margin changes are conditional screens, not
probabilistic adequacy findings.}
\label{fig:ercot}
\end{figure}

In the matched normal 2030 cases, grid-dominant and not-yet-commissioned
firm-low-carbon portfolios leave 4.145 GW on the grid. Firm onsite supply has
3.316 GW accredited onsite capacity and leaves 0.829 GW; the
renewable--storage--grid portfolio has 2.777 GW accredited and leaves 1.368 GW.
Figure~\ref{fig:portfolio} preserves LOLE and expected unserved energy as null.
The screening cannot identify which portfolio reduces probabilistic adequacy
risk, so H2 and H3 remain indeterminate.

\begin{figure}[htbp]
\centering
\includegraphics[width=\linewidth]{figure_07.pdf}
\caption{Accredited-capacity and reserve-margin replacement for the
prespecified adequacy figure. LOLE and expected unserved energy are unavailable
and are not inferred.}
\label{fig:portfolio}
\end{figure}

\subsection{National annual-electricity context}
\label{sec:results-national}

The accelerated case equals 2.636\% of the EIA baseline US annual electricity
in 2037 and 2.255\% in 2050 (Figure~\ref{fig:national}). The decline after full
target is caused by growth in the national denominator, not lower facility
demand. These shares contextualize scale; they are neither national security
thresholds nor evidence of regional feasibility. National peak share, fuel
demand and operational emissions remain null because the evidence contract did
not identify consistent outputs for them.

\begin{figure}[htbp]
\centering
\includegraphics[width=\linewidth]{figure_08.pdf}
\caption{Conditional facility electricity as a share of EIA baseline and
high-electricity-demand annual totals. National peak, fuel and emissions
results are omitted rather than invented.}
\label{fig:national}
\end{figure}

\subsection{Gate outcomes and realism classification}
\label{sec:results-gates}

Across 20,720 active-case gate evaluations, 8,724 pass, 2,228 fail under
explicit schedule, supply, thermal or water stresses, and 9,768 are
indeterminate. Manufacturing throughput, evidence sufficiency and regional
adequacy are indeterminate for every active case. Thermal and water gates are
indeterminate in ordinary cases and fail only in the declared 80\%-capacity
stress. Schedule precedence and supply deliverability account for 818 failed
evaluations. Consequently all 336 nonzero trajectories are indeterminate and
no first joint-pass year exists (Figure~\ref{fig:gates}). H1--H4 all remain
indeterminate under their predeclared definitions.

\begin{figure*}[htbp]
\centering
\includegraphics[width=0.94\textwidth]{figure_09.pdf}
\caption{Trajectory classification and gate-status map. Explicit failures and
missing-evidence indeterminacy are distinct; neither is relabelled for a simpler
overall verdict.}
\label{fig:gates}
\end{figure*}

\subsection{Uncertainty, sensitivity and convergence}
\label{sec:results-uncertainty}
\label{sec:results-convergence}

For the full target under the baseline correlated design and 32,768
confirmation samples, median wafer starts are 50.445 million/year and median
facility electricity is 125.107 TWh/year. Marginal p95 facility peak is
91.542 GW and marginal p95 withdrawal is 1.571 billion m$^3$/year. These p95
values are separate marginal quantiles and must not be combined into a joint
design point. Feasibility probability is null because indispensable project
capacity and probabilistic adequacy inputs are absent.

The worst 16,384-to-32,768 relative change is 0.289999\%, below the locked 1\%
criterion across independence, baseline dependence and both correlation
stresses (Table~\ref{tab:verification}). For the independent conditional-peak
diagnostic, logic dies per module has total-order index 0.5397, facility
electricity intensity 0.2404, die area 0.1495, module rated power 0.0371 and
effective yield 0.0324. The top two contribute 78.0212\% of normalized
total-order mass, below the predeclared 80\% threshold; no favorable rounding
is applied. The dependence-aware surrogate has $R^2=0.9405$ and retains logic
dies per module, electricity intensity and die area as the leading three.
Figure~\ref{fig:sensitivity} reports intervals and convergence.

\begin{figure*}[htbp]
\centering
\includegraphics[width=0.94\textwidth]{figure_10.pdf}
\caption{Independent Sobol and dependence-aware Shapley diagnostics with
bootstrap intervals, plus sample-doubling convergence. The outcome is
conditional facility peak, not the unavailable joint-feasibility outcome.}
\label{fig:sensitivity}
\end{figure*}

\section{Discussion}

\subsection{What the public evidence supports}

The central scientific answer has two parts. First, the currently admitted
public record does not demonstrate that the announced project is realistic.
The conclusion follows from missing target semantics and missing phase
manufacturing, electricity, thermal, water, interconnection and probabilistic
adequacy capacity---not from a finding that a physical law is violated. The
proper classification is therefore indeterminate, not infeasible.

Second, the primary admitted interpretation is energy-security-material. A
central 16.58-GW industrial peak would be too large to treat as an ordinary
unmodeled load addition, and the broad epistemic uncertainty produces an even
larger tail. The calculation establishes a planning-relevant order of
magnitude under transparent assumptions. It does not establish an official
load, an interconnection obligation, or a regulatory result.

The distinction matters for both favorable and unfavorable narratives. A
positive reserve-margin change under one capacity counterfactual is not proof
of improved reliability; neither is an explicit failure in a deliberately
stressed water or schedule case proof that the ordinary project fails. By
retaining tri-state gates and separate counterfactuals, the model exposes where
the conclusion depends on evidence rather than replacing missing evidence with
an analyst's preference.

\subsection{Energy-conversion and management contribution}

The production-to-energy chain adds four methodological safeguards relevant to
other emerging industrial loads. It prevents dimensional collapse between
product output and factory input; maintains a closed facility electricity and
heat allocation; types water intensities to avoid double-counting reuse; and
separates reserve-margin screening from probabilistic adequacy. The same
architecture can be applied to other public industrial claims whose physical
scale is clearer than their utility implications.

The analysis also identifies where additional evidence has highest decision
value. A dimensional product definition and phase wafer capacity would test the
manufacturing translation. Phase electrical, thermal and water designs would
bound facility gates. An executable point and size of interconnection,
portfolio-specific accredited capacity and hourly probabilistic studies would
permit regional adequacy assessment. These observations identify scientific
measurements and disclosures; they are not procurement, permitting or
infrastructure-planning recommendations.

\subsection{Limitations}

The analysis is prospective and comparator-based, not an empirically validated
digital twin. No independent numerical Terafab holdouts exist, and the public
benchmarks used to define ranges cannot be reused as validation observations.
The continuous die geometry is a screening approximation; actual layouts,
node mix, memory, packaging, test and yield learning are unknown. Total-facility
intensity compresses process heterogeneity, while the heat result lacks
temperature-grade and sink information. Water availability and discharge are
not site models.

Regional screening omits power flow, congestion, unit commitment, outage and
weather chronology, load flexibility, fuel networks, price formation, and
probabilistic LOLE/EUE. ERCOT counterfactuals stop at 2030. National results are
annual-energy shares only and cannot imply regional or hourly security. The
uncertainty distributions are epistemic scenarios, not frequencies of future
project outcomes; their marginal p95 values are not a joint worst case.
Construction embodied energy, supply-chain benefits, costs, emissions, project
finance and macroeconomic effects are outside scope.

\section{Conclusions}

Under a transparent interpretation of one terawatt of rated compute-device
power manufactured annually, the announced output implies a central
requirement of 68.29 million 300-mm wafer starts/year, 135.17 TWh/year,
16.58 GW coincident electricity and heat rejection, and 568.72 million
m$^3$/year external water withdrawal. Those requirements are large enough to
be material to ERCOT and US electricity planning. Yet all nonzero trajectories
remain evidentially indeterminate: current public information does not define
the target or demonstrate the manufacturing, utility and adequacy capacities
needed to realize it. The defensible answer is therefore that project realism
is not presently demonstrated, not that the project is proven infeasible.

This conclusion is deliberately falsifiable. Public disclosure of dimensional
output, phase capacity, facility resource designs, interconnection and
probabilistic adequacy results would convert the unknown gates into tests. Until
then, the reproducible model supplies conditional resource implications while
preserving the boundary between analytical estimates and verified project
facts.

\section*{Data and code availability}

The source-available code, frozen public-evidence extracts, scenario contracts,
notebook and deterministic regeneration instructions are provided at
\url{https://github.com/KNOWDYN/terafab-decision-twin}. Noncommercial academic
use is governed by \code{ACADEMIC_LICENSE.md}; commercial use requires a
separate written license under \code{COMMERCIAL_LICENSE.md}. The final archival
release, Git commit and DOI must be inserted after the Step 7 repository
publication and archival deposit. Generated outputs are analytical estimates,
not verified Terafab operating facts.

\section*{CRediT authorship contribution statement}

Khalid M. Saqr: Conceptualization, Methodology, Software, Validation, Formal
analysis, Investigation, Data curation, Writing---original draft,
Writing---review and editing, Visualization, Supervision and Project
administration.

\section*{Funding}

The funding statement is author-controlled and must be supplied accurately
before journal submission; this reproducibility package does not infer a
funding source.

\section*{Declaration of competing interest}

The author is Managing Director of KNOWDYN Ltd, which owns the source-available
software used in this work. The study is an independent KNOWDYN analysis and is
not affiliated with, endorsed by, authorized by, sponsored by, or connected to
Terafab or its employees. References to external organizations are analytical
and do not imply access to verified internal data or endorsement.

\section*{\raggedright Declaration of generative AI and AI-assisted technologies}

During preparation of this work, the author used OpenAI Codex (GPT-5 family,
accessed 1 August 2026) to assist source-code implementation, automated testing,
reproducible data-visualization code, manuscript organization and language
drafting. All figures are deterministic Matplotlib renderings of the cited
machine-readable outputs; no generative imagery or fabricated data were used.
After using this service, the author reviewed and edited the content as needed
and takes full responsibility for the content of the publication. The exact
service model/version must be confirmed against the author's access record
before submission.

\bibliographystyle{unsrt}
\bibliography{references}

\end{document}
